\documentclass[11pt]{article}
\usepackage[a4paper,margin=1in]{geometry}
\usepackage{amsmath,amsthm,amssymb}
\usepackage{mathtools}

\newtheorem{theorem}{Theorem}[section]
\newtheorem{proposition}[theorem]{Proposition}
\newtheorem{lemma}[theorem]{Lemma}
\newtheorem{corollary}[theorem]{Corollary}
\theoremstyle{definition}
\newtheorem{definition}[theorem]{Definition}
\newtheorem{remark}[theorem]{Remark}

\newcommand{\R}{\mathbb{R}}
\newcommand{\Fr}{\mathcal{A}}
\newcommand{\Om}{\Omega}
\newcommand{\eps}{\varepsilon}
\newcommand{\rstar}{\rho_{*}}
\newcommand{\rhad}{\widehat\rho}
\newcommand{\X}{\mathcal X}

\title{Integrated Action, Kinetic Bound,\\
       and the Weighted Metric Trace Lemma\\
       (Route~$\delta$, Plan~2, Building Block \textsc{WeightedMetricTrace})}
\author{}
\date{}

\begin{document}
\maketitle

\begin{abstract}
This building block chains three ingredients into a pointwise metric
distance bound at a radius $\rhad$ near $\rho_\delta=1-\kappa\sqrt{\delta_T}$.
The starting point is the integrated quadratic Fusco--Julin oscillation
$(B^2\text{-int})$ of \cite{SpaceTimeIdentity}, which replaces any
attempt at a pointwise $L^2$ radial trace by its $d\mu$-averaged
form. Combining $(B^2\text{-int})$ with the per-$\rho$ Cauchy deficit
bound of \cite{MetricFramework} yields the integrated action bound
$\int(d_\X(F_\rho,B_1)^2+|\dot F_\rho|_\X^2)\,d\mu \le C\,\delta_T$.
A Talenti / bad-set decomposition then divides by $-t'(\rho)$ on the
good set and uses the elementary uniform $L^1$-Lipschitz bound (A0)
on the bad set, giving the kinetic bound (K) in Lebesgue form
$\int|\dot F_\rho|_\X^2\,d\rho\le C'\delta_T$. The abstract weighted
metric trace lemma then produces $\rhad\in[\rho_\delta-C_*\delta,\rho_\delta]$
with $d_\X(F_{\rhad},B_1)^2\le C_*\,\delta_T$. The chain uses no
pointwise (R)/(1.2) anywhere.
\end{abstract}

\section{Statements}\label{sec:statements}

Fix $n\ge 2$, $R\ge 2$, $\rstar\in(0,1)$, $\kappa>0$. Let $\Om\subset B_R$
open with $|\Om|=\omega_n$, $u=u_\Om$ the variational torsion function,
$\delta:=\delta_T(\Om)\le\delta_0(n,R,\rstar)$. Set $E_\rho:=\{u>t(\rho)\}$
with $|E_\rho|=\omega_n\rho^n$, $\rho_\delta:=1-\kappa\sqrt\delta$,
$\mu(d\rho):=(-t'(\rho))d\rho$. Define
$F_\rho:=\rho^{-1}E_\rho\subset B_{R/\rstar}$, $\gamma(\rho):=[F_\rho]\in\X$.

\begin{theorem}[Integrated action bound]\label{thm:52int}
Conditional on (F) of \cite{CentroidBound} and $(B^2\text{-int})$ of
\cite{SpaceTimeIdentity}:
\begin{equation}\label{eq:52int}
   \int_{\rstar}^{\rho_\delta}\bigl(d_\X(F_\rho,B_1)^2+|\dot F_\rho|_\X^2\bigr)\,d\mu(\rho)
   \le C(n,R,\rstar,\kappa)\,\delta_T(\Om).
\end{equation}
\end{theorem}

\begin{theorem}[Kinetic bound (K), Lebesgue form]\label{thm:K}
\begin{equation}\label{eq:K}
   \int_{\rstar}^{\rho_\delta}|\dot F_\rho|_\X^2\,d\rho
   \le C'(n,R,\rstar,\kappa)\,\delta_T(\Om).
\end{equation}
\end{theorem}

\begin{theorem}[Weighted metric trace]\label{thm:trace}
There exist $C_*$ and $\rhad\in[\rho_\delta-C_*\delta,\rho_\delta]$ such that
$d_\X(F_{\rhad},B_1)^2\le C_*\,\delta_T(\Om)$.
\end{theorem}

\section{The integrated action bound (5.2-int)}\label{sec:52int}

From \cite{MetricFramework} Theorem (T),
$|\dot F_\rho|_\X\le C\inf_a\int|V_\rho-H_{z_\rho,\rho}-a\cdot\nu_\rho|$;
squaring gives
$|\dot F_\rho|_\X^2\le 2[(\int|V_\rho-\bar V_\rho|)^2+(\int|H_{z_\rho,\rho}-\bar V_\rho|)^2]$.

\paragraph{$D_H$-piece.} The unconditional Cauchy--Schwarz on $V_\rho$
(algebraic from \cite{LevelSetIdentity}):
$(\int|V-\bar V|)^2\le C\, D_H$.
Integrating: $\int(\int|V-\bar V|)^2 d\mu\le C\int D_H\,d\mu\le C\delta$.

\paragraph{Radial-trace piece.} Using $|H-\bar V|\le||x-z|/\rho-1|+|1-\bar V|$
plus Cauchy--Schwarz on $\partial^*E_\rho$:
$\int(\int|H-\bar V|)^2 d\mu\le C\delta$ via $(B^2\text{-int})$ of
\cite{SpaceTimeIdentity}, plus perimeter excess integration and the
slice-rearrangement $|E\Delta B_\rho|^2\le C\Pi$ + (G3).

\paragraph{Distance piece.} $d_\X(F_\rho,B_1)\le\rho^{-n}|E_\rho\Delta B_\rho(z_\rho)|$
pointwise, so $\int d_\X^2 d\mu\le C\int|E_\rho\Delta B_\rho|^2 d\mu\le C\delta$
via the slice-rearrangement.

Combining gives \eqref{eq:52int}. \hfill\qed

\begin{remark}
The (T1²-int) input from \cite{SpaceTimeIdentity} replaces the conditional
pointwise (R)/(1.2). The chain is rate-preserving thanks to the
slicing-then-squaring trick.
\end{remark}

\section{The kinetic bound (K)}\label{sec:K}

\begin{lemma}[Uniform metric Lipschitz, (A0)]\label{lem:A0}
$|\dot F_\rho|_\X\le L_0(n,R,\rstar)$ a.e.
\end{lemma}

\begin{proof}
Nestedness gives $|E_{\rho'}\Delta E_\rho|\le n\omega_n|\rho'-\rho|$.
Split $F_\rho\Delta F_{\rho'}$ through $(\rho')^{-1}E_\rho$ and use
the dilation estimate to get $d_\X([F_{\rho'}],[F_\rho])\le L_0|\rho'-\rho|$.
\end{proof}

\begin{lemma}[Bad-set Lemma]\label{lem:badset}
$|B_{\tau_0}|:=|\{\rho:-t'(\rho)<\tau_0=\rstar/(2n)\}|\le K_2\delta$.
\end{lemma}

\begin{proof}
On $B_{\tau_0}$, $|\psi(\rho)|=-t_B'-(-t')\ge\rstar/(2n)$. By the
integrated Talenti gap $\int|\psi|d\rho\le K_1\delta$,
$(\rstar/(2n))|B_{\tau_0}|\le K_1\delta$.
\end{proof}

\begin{proof}[Proof of Theorem~\ref{thm:K}]
Split $[\rstar,\rho_\delta]=G\sqcup B_{\tau_0}$ with $G$ the good set
where $-t'\ge\tau_0$. On $G$:
$\int_G|\dot F_\rho|^2 d\rho\le\tau_0^{-1}\int_G|\dot F_\rho|^2(-t')d\rho
\le(2n/\rstar)\cdot C\delta$. On $B_{\tau_0}$: by (A0),
$\int_{B_{\tau_0}}|\dot F_\rho|^2 d\rho\le L_0^2 K_2\delta$.
Adding gives \eqref{eq:K}.
\end{proof}

\section{The weighted metric trace lemma}\label{sec:agent4}

\begin{theorem}[Abstract weighted metric trace]\label{thm:abstract-trace}
Assume (H1) $\int(d_\X^2+|\dot\gamma|^2)d\mu\le C\delta$, (H2)
$\mu([a,b])\ge c(b-a)-C\delta$, (H3) $\mu_\rho\le M$, (5.14)
$\int|\dot\gamma|^2 d\rho\le C'\delta$. Then there is
$\rhad\in[\rho_\delta-C_*\delta,\rho_\delta]$ with
$d_\X(\gamma(\rhad),B_1)^2\le C_*\delta$.
\end{theorem}

\begin{proof}
\emph{Step 1: macroscopic Markov on a fixed sub-interval.} Fix
$L^*:=(\rho_\delta-\rstar)/4$, $J^*:=[\rho_\delta-L^*,\rho_\delta-L^*/2]$.
By (H2), $\mu(J^*)\ge cL^*/4$. By (H1) + Markov on $J^*$, there is
$\rho_0\in J^*$ with $d_\X(\gamma(\rho_0),B_1)^2\le K_1\delta$ where
$K_1=4C/(cL^*)$.

\emph{Step 2: kinetic transport.} For any $\rhad\in I_\delta:=[\rho_\delta-K_2\delta,\rho_\delta]$,
AGS Thm 1.1.2 + Cauchy--Schwarz:
$d_\X(\gamma(\rho_0),\gamma(\rhad))^2\le(\rhad-\rho_0)\int|\dot\gamma|^2 d\rho
\le L^*\cdot C'\delta$ by (5.14).

\emph{Step 3: triangle.} $d_\X(\gamma(\rhad),B_1)^2\le 2K_1\delta+2L^*C'\delta=:C_*\delta$.
\end{proof}

\paragraph{Verification of hypotheses for the torsion application:}
\begin{itemize}
\item (H1) = Theorem~\ref{thm:52int}.
\item (H2) interval bound from Talenti / Nicola--Tilli profile gap
\cite{LevelSetIdentity}.
\item (H3) $-t'\le 1/n$ via Talenti.
\item (5.14) = Theorem~\ref{thm:K}.
\end{itemize}

\begin{proof}[Proof of Theorem~\ref{thm:trace}]
Apply Theorem~\ref{thm:abstract-trace} to $\gamma=[F_\rho]$.
\end{proof}

\section{Conditionality}\label{sec:cond}

\begin{itemize}
\item Theorem~\ref{thm:52int}: consumes $(B^2\text{-int})$ + (G3) of
\cite{SpaceTimeIdentity}, unconditional given (F) of \cite{CentroidBound}.
\item Theorem~\ref{thm:K}: plus elementary (A0) and bad-set Lemma,
both unconditional.
\item Theorem~\ref{thm:trace}: plus (H2), (H3) from Talenti, unconditional.
\end{itemize}

The whole block is unconditional given (F). Pointwise (R)/(1.2) is
never invoked.

\begin{thebibliography}{99}

\bibitem{LevelSetIdentity}
Building block \textsc{LevelSetIdentity}.

\bibitem{MetricFramework}
Building block \textsc{MetricFramework}.

\bibitem{CentroidBound}
Building block \textsc{CentroidBound}.

\bibitem{SpaceTimeIdentity}
Building block \textsc{SpaceTimeIdentity}.

\bibitem{BoundaryLayerTransfer}
Building block \textsc{BoundaryLayerTransfer}.

\bibitem{AGS2008}
L.~Ambrosio, N.~Gigli, G.~Savar\'e, \emph{Gradient Flows in Metric
Spaces}, Birkh\"auser, 2008. Thm 1.1.2.

\bibitem{Talenti1976}
G.~Talenti, Ann. Mat. Pura Appl. (4) \textbf{110} (1976), 353--372.

\bibitem{Maggi2012}
F.~Maggi, Cambridge University Press, 2012.

\end{thebibliography}

\end{document}
